\documentclass[../Main.tex]{subfiles}

\begin{document}
\section{Definitions of Chaos}
Chaos generally involves:
\begin{itemize}
    \item arbitrarily close orbits eventually separate;
    \item even the smallest subsets mix (there are no isolated orbits).
\end{itemize}
It is impossible to make long-time predictions for chaotic systems, because any small error in the initial measurement will be propagated into a much larger error in the eventual solution.

\begin{definition}{Sensitive dependence on initial conditions}
    Consider a continuous map $F : I \mapsto I$ of a bounded interval $I \subset R$ onto itself. Let $\Lambda$ be a set invariant under $F$. $F$ has \underline{sensitive dependence on}\\\underline{initial conditions} (SDIC) on $\Lambda$ if, for any $\epsilon > 0$ and any $x \in \Lambda$, there exists $y \in \Lambda, \delta > 0$ and $n \in \N$ such that $\abs{x - y} < \epsilon$ and $\abs{F^n(x) - F^n(y)} > \delta$.
\end{definition} %TODO: check -- I wonder if we can choose any $\delta$ and then it is $n(\delta)$.
\begin{definition}{Topological transitivity}
    With $F$ as above defined, it is \underline{topologically transitive} (TT) on $\Lambda$ if for any non-empty open sets $U, V \subset \Lambda$, there exists $n > 0$ such that $F^n(U) \cap V \neq \emptyset$.
\end{definition}
\begin{remark}
    This definition is most interesting if $U$ and $V$ do not intersect, then we are saying that orbits starting in $U$ eventually reach $V$. In other words, $\Lambda$ cannot be decomposed into disjoint, invariant subsets.
\end{remark}
\begin{example}[TT but not SDIC]
    Consider the rotation map on $[0, 1)$ given by:
    \begin{equation*}
        x_{n+1} = x_n + \omega \mod 1
    \end{equation*}
    If $\omega$ is irrational, the orbits are dense (TT) but the separation of two close points stays the same (not SDIC).
\end{example}
\begin{example}[SDIC but not TT]
    Consider a map given by:
    \begin{equation*}
        x_{n+1} =
        \begin{cases}
            2(-1-x_n) & x_n < -\frac12 \\
            2x_n & -\frac12 \leq x_n \leq \frac12 \\
            2(1-x_n) & x_n > \frac12
        \end{cases}
    \end{equation*}
    This is not a topologically transitive map because the subsets $x > 0$ and $x < 0$ are independently invariant. However, it is SDIC since $\abs{F_x} = 2 > 1$. Consider perturbing $x_n \mapsto x_n + \delta_n$.
    \begin{equation*}
        x_{n+1} + \delta_{n+1} = F(x_n + \delta_n) = F(x_n) + F_x(x_n) \delta_n
    \end{equation*}
    or, $\delta_{n+1} = F_x(x_n) \delta_n$, and so this is unstable.
\end{example}
\begin{definition}{D-chaos}
    A function $F : \Lambda \mapsto \Lambda$ is \underline{D-chaotic} if:
    \begin{enumerate}
        \item $F$ has sensitive dependence to initial conditions on $\Lambda$;
        \item $F$ is topologically transitive on $\Lambda$;
        \item periodic points are dense in $\Lambda$.
    \end{enumerate}
\end{definition}
\end{document}